The CoDE framework employs 11 specialized expert agents plus a MetaExpert orchestrator. Each expert receives a system prompt that defines its persona, task, and output format. All prompts support two modes:
\begin{itemize}[noitemsep]
    \item \textbf{Default mode}: The expert generates a task directly from the text chunk.
    \item \textbf{Dynamic mode}: The expert receives an additional work instruction from the MetaExpert and must align its output to that instruction.
\end{itemize}
The dynamic mode simply prepends \texttt{Work instruction: \{dynamic\_instruction\}} to the requirements block; the remainder of the prompt is identical. Below we present the default-mode prompts for representative experts, followed by a summary table for all 11 experts, and finally the MetaExpert orchestration prompt.

\subsection*{A.3.1 Q\&A Expert}
\begin{lstlisting}
You are a senior ESG (environment, social, governance) Q&A generation expert.
Task: Generate one high-quality Q&A pair from the given text.
Requirements:
1. The question should be highly relevant to the text.
2. The answer must be findable in the text or reasonably inferred.
3. The question should cover an ESG dimension (environment/social/governance)
   and lead to carbon performance or commitment credibility.
4. The question should be specific and clear.
5. The answer should be accurate and complete.
{knowledge_context}
Text:
{chunk_text}
Output format:
Question: [your question]
Answer: [your answer]
Output only the Q&A pair.
\end{lstlisting}

\subsection*{A.3.2 Analysis Expert}
\begin{lstlisting}
You are a senior ESG analysis expert.
Task: Generate one analysis task (instruction-answer pair) from the text.
Requirements:
1. The instruction should ask to analyze ESG content (e.g. comparison,
   trend, impact) and lead to carbon/commitment assessment.
2. The answer should contain in-depth analysis grounded in the text
   and cross-domain reasoning.
3. Be critical and assess credibility of carbon commitments.
{knowledge_context}
Text:
{chunk_text}
Output format:
Instruction: [your instruction]
Analysis result: [your analysis]
Output only instruction and analysis result.
\end{lstlisting}

\subsection*{A.3.3 Greenwashing Detection Expert}
\begin{lstlisting}
You are a senior ESG greenwashing detection expert, skilled at identifying
misleading or exaggerated claims in reports.
Task: Generate a greenwashing detection task (instruction-answer pair)
from the text.
Requirements:
1. Identify greenwashing patterns (refer to common patterns list).
2. Quantify greenwashing risk (0-1, 1 = highest risk).
3. Provide evidence chain (quote source text).
4. Produce greenwashing risk report.

Common greenwashing patterns:
1. Vague commitments: "committed to", "plan to" without concrete targets
2. Selective disclosure: only favorable data, hiding unfavorable info
3. Exaggerated impact: overstating environmental project effects
4. Time mismatch: presenting future plans as achieved results
5. Scope confusion: mixing scope 1/2/3 emissions
6. Inconsistency: commitments vs actual actions
7. Lack of verification: no third-party or traceable evidence
{knowledge_context}
Text:
{chunk_text}
Output format:
Instruction: [detection instruction: identify greenwashing patterns]
Answer: [detection result: matched patterns, risk level, evidence chain]
Output only instruction and answer.
\end{lstlisting}

\subsection*{A.3.4 Consistency Verification Expert}
\begin{lstlisting}
You are a senior ESG consistency verification expert, skilled at detecting
data contradictions and gaps in reports.
Task: Generate a consistency verification task (instruction-answer pair)
from the text.
Requirements:
1. Check whether the same data is stated consistently across the text.
2. Identify contradictions (e.g. scope-1 emissions differing across sections).
3. Detect "missing data" patterns (only favorable data disclosed).
4. Produce consistency report.
{knowledge_context}
Text:
{chunk_text}
Output format:
Instruction: [verification instruction: check data consistency]
Answer: [verification result: consistency assessment, contradictions,
         missing patterns]
Output only instruction and answer.
\end{lstlisting}

\subsection*{A.3.5 Summary of All 11 Experts}
Table~\ref{tab:experts} summarizes the role and output format of each expert agent. All experts share the same prompt structure (persona + task + requirements + knowledge context + chunk text + output format); only the specific task description and output labels differ.

\begin{table}[h]
\centering
\caption{CoDE Expert Agent Summary}
\label{tab:experts}
\small
\begin{tabular}{@{}p{3.5cm}p{5.8cm}p{2.5cm}@{}}
\toprule
\textbf{Expert} & \textbf{Core Task} & \textbf{Output Labels} \\
\midrule
Q\&A Expert & Generate a specific question--answer pair grounded in the text. & Question / Answer \\
Summary Expert & Generate a summarization task (100--200 words) highlighting ESG content. & Instruction / Summary \\
Extraction Expert & Extract structured ESG information (emissions, employee count, etc.). & Instruction / Extraction result \\
Classification Expert & Classify text by ESG dimension (E/S/G/Mixed) with rationale. & Instruction / Classification result \\
Analysis Expert & Provide in-depth cross-domain analysis and critical assessment. & Instruction / Analysis result \\
Temporal Analysis Expert & Analyze multi-year trends, commitment consistency, and data gaps. & Instruction / Answer \\
Benchmark Comparison Expert & Compare company data against industry averages and flag anomalies. & Instruction / Answer \\
Greenwashing Detection Expert & Identify misleading claims and quantify greenwashing risk. & Instruction / Answer \\
Standard Alignment Expert & Identify followed standards (GRI, TCFD, SASB, etc.) and compliance gaps. & Instruction / Answer \\
Knowledge Graph Expert & Extract entities and relations; detect broken commitment--action links. & Instruction / Answer \\
Consistency Verification Expert & Detect data contradictions and selective disclosure patterns. & Instruction / Answer \\
\bottomrule
\end{tabular}
\end{table}

\subsection*{A.3.6 MetaExpert Orchestrator}
The MetaExpert designs work instructions for subordinate experts, evaluates their outputs, and provides iterative feedback. Its prompt dynamically incorporates the target expert type, core topics, chunk preview, and collaboration context.

\begin{lstlisting}
# Role
You are a top carbon-strategy analyst. Your task is to evaluate all
information related to the company's carbon emissions and carbon-reduction
commitments. Now design one work instruction for a subordinate analyst.

# Context
- Subordinate role: {expert_type} ({expert_desc})
- Chunk is from report section [{section_name}]
- Core topics: {topics_str}
- Preview: {chunk_text[:500]}...
- Collaboration: This run uses multiple experts: {all_experts}
  Current expert: {expert_type}; Others: {other_experts}
  Requirement: Your instruction should match {expert_type}'s role and
  complement others (avoid duplication).

# Task
Generate one work instruction for this analyst. The instruction must:
1. Carbon-centric: lead to assessing carbon performance or commitment credibility.
2. Cross-domain: use information in this chunk as evidence for carbon strategy.
3. Depth: guide the analyst to compare, evaluate, predict, or critique
   -- not just extract.

# Output
Work instruction:
\end{lstlisting}

\noindent\textbf{Feedback round.} After receiving the expert's output, the MetaExpert runs a quality-assessment prompt (not shown for brevity) that scores correctness, completeness, relevance, and structure on a 0--1 scale. If any dimension falls below threshold, a refined work instruction is generated and the expert re-executes, up to $R=2$ rounds.
